% Article on the transient absorption of LiF from the non-linear susceptibilities (plan: Outline.md).
% All the text is in this file (no separate files for the sections).
\documentclass[aps,prb,reprint,superscriptaddress,amsmath,amssymb]{revtex4-2}

\usepackage[utf8]{inputenc}
\usepackage[english]{babel}
\usepackage{bm}
\usepackage{graphicx}
\usepackage{booktabs}
\usepackage[colorlinks=true,linkcolor=blue,citecolor=blue,urlcolor=blue]{hyperref}

% Notation of the article: the same macros as the technical notes (Notes/header.tex), so that the two documents use the
% same symbols. Add here the macros of every recurring symbol.
\newcommand{\chin}[2]{\chi_{(#1,#2)}}          % MPPI key chi_(n,m)
\newcommand{\dchi}{\delta\chi}                  % pump induced change of chi
\newcommand{\deps}{\Delta\varepsilon}           % transient Delta eps = 4 pi dchi
\newcommand{\EP}{E_{\mathrm{P}}}                % pump field
\newcommand{\Ep}{E_{\mathrm{p}}}                % probe field
\newcommand{\wP}{\omega_{\mathrm{P}}}           % pump frequency
\newcommand{\tP}{t_{\mathrm{P}}}                % pump maximum
\newcommand{\tp}{t_{\mathrm{p}}}                % probe time
\newcommand{\fenv}{f}                           % pump envelope
\newcommand{\chineq}{\chi^{\mathrm{neq}}}       % non-equilibrium response
\newcommand{\PPp}{P_{\mathrm{Pp}}}              % polarization with pump and probe
\newcommand{\PPo}{P_{\mathrm{P}}}               % polarization with the pump only
\newcommand{\EOP}{E_{0\mathrm{P}}}              % peak amplitude of the pump
\newcommand{\chiz}{\chi^{(3)}_{0}}              % 6 chi3(w;w,wP,-wP)
\newcommand{\chip}{\chi^{(3)}_{+}}              % 3 chi3(w;w-2wP,wP,wP)
\newcommand{\chim}{\chi^{(3)}_{-}}              % 3 chi3(w;w+2wP,-wP,-wP)


\begin{document}

\title{Understanding pump and probe signal from excitonic non linear response functions}

% Authors and affiliations: placeholders, to be completed
\author{Marco D'Alessandro}
\affiliation{[affiliation]}
\author{Claudio Attaccalite}
\affiliation{[affiliation]}
\author{Davide Sangalli}
\affiliation{[affiliation]}

\date{\today}

\begin{abstract}
[Abstract to be written.]
\end{abstract}

\maketitle

% ====================================================================================================
\section{Introduction}
\label{sec:intro}

% To be written by Davide (Outline.md, Sec. 1).
\textit{[Among the points: motivation of the approach. RT simulations with explicit pump
and probe fields give the transient response, but they do not explain the origin of its features; the
non-linear susceptibilities are properties of the system and the pump enters the response in a way fixed by the
formalism (Sec.~\ref{sec:theory}), so that the features of the signal can be traced back to them.]}

% ====================================================================================================
% Provisional title
\section{Transient response: real-time simulations and non-linear susceptibilities}
\label{sec:theory}

% Plan: Outline.md, Sec. 2 (general: no LiF, no codes). The motivation of the approach is in the Introduction.

\subsection{Pump-probe response}
\label{sec:theory_neq}

In a pump-probe setup the system is driven by a pump field $\bm{E}_{\mathrm{P}}$ and by a weak probe field
$\bm{E}_{\mathrm{p}}$, delayed with respect to the pump. We take the two fields along the same direction, and denote by
$\EP$, $\Ep$ and $P$ the components of the fields and of the polarization along it. The transient response is the
polarization induced by the probe in the presence of the pump, divided by the probe field,
\begin{equation}
\chineq[\EP](\omega) = \frac{\PPp(\omega)-\PPo(\omega)}{\Ep(\omega)} ,
\label{eq:chineq}
\end{equation}
where $\PPp$ is the polarization with both the fields and $\PPo$ the one with the pump only. In a real-time (RT)
simulation $\PPp(t)$ and $\PPo(t)$ are obtained from the time evolution of the system in the fields, and
$\chineq$ from their Fourier transforms; the pump-only simulation does not depend on the probe and is the same for
all the delays. Eq.~\eqref{eq:chineq} contains the response to all orders in the pump field, and it is linear in the
probe field as long as the probe is weak.
% TODO: references for the RT pump-probe simulations

We consider an impulsive probe at the time $\tp$,
\begin{equation}
\Ep(t) = E_{0\mathrm{p}}\,\delta(t-\tp) , \qquad
\Ep(\omega) = \frac{E_{0\mathrm{p}}}{2\pi}\,e^{i\omega\tp} ,
\label{eq:delta_probe}
\end{equation}
which models a probe pulse much shorter than the period of the pump and than the time scale of the response. All the
probe frequencies have the same amplitude and a phase fixed by $\tp$.

With this probe the transient response depends on the probe time, i.e.\ on the delay between probe and pump. Its
change induced by the pump,
\begin{equation}
\dchi(\omega,\tp) = \chineq(\omega,\tp)-\chi^{(1)}(\omega) ,
\label{eq:dchi}
\end{equation}
where $\chi^{(1)}$ is the linear susceptibility at equilibrium, is the quantity that we express in terms of the
equilibrium susceptibilities in the following. It gives the change of the dielectric function with respect to the
equilibrium one, $1+4\pi\chi^{(1)}(\omega)$,
\begin{equation}
\deps(\omega,\tp) = 4\pi\,\dchi(\omega,\tp) ,
\label{eq:deps}
\end{equation}
whose imaginary part is the change of the absorption.

\subsection{Expansion in the equilibrium susceptibilities}
\label{sec:theory_expansion}

To express $\dchi$ in terms of the equilibrium susceptibilities we write fields and polarization as superpositions of
complex exponentials~\cite{Boyd},
\begin{equation*}
\bm{E}(t) = \int d\omega\,\bm{E}(\omega)\,e^{-i\omega t} , \qquad
\bm{P}(t) = \int d\omega\,\bm{P}(\omega)\,e^{-i\omega t} ,
\end{equation*}
so that $\bm{E}(\omega) = (2\pi)^{-1}\int dt\,\bm{E}(t)\,e^{i\omega t}$ and $\bm{E}(-\omega) = \bm{E}(\omega)^*$ for a
real field (atomic units are used throughout). The polarization is expanded in powers of the field,
$\bm{P} = \bm{P}^{(1)}+\bm{P}^{(2)}+\bm{P}^{(3)}+\dots$, and the susceptibility of order $n$ is defined by
\begin{align}
P^{(n)}_i(\omega) &= \int d\omega_1\dots d\omega_n\,\delta\Big(\omega-\sum_{k=1}^n\omega_k\Big) \nonumber\\
&\quad\times\chi^{(n)}_{ij_1\dots j_n}(\omega;\omega_1,\dots,\omega_n)\prod_{k=1}^n E_{j_k}(\omega_k) ,
\label{eq:Pn}
\end{align}
where the repeated Cartesian indices are summed over. The susceptibilities $\chi^{(n)}$ are properties of the system
at equilibrium. Their first frequency argument is the output frequency, equal to the sum of the input ones, and they
can be taken symmetric under the simultaneous permutation of the input frequencies and of the corresponding Cartesian
indices (intrinsic permutation symmetry)~\cite{Boyd}. With the fields along the same direction only the diagonal
component along it enters, which we denote by $\chi^{(n)}$.

We insert the sum of the two fields, $E = \Ep+\EP$, in Eq.~\eqref{eq:Pn}. The terms with the pump only give
$\PPo$, which cancels in Eq.~\eqref{eq:chineq}, and we keep the terms linear in the probe. In $\chi^{(n)}$ the probe
can take any of the $n$ input places, which give the same contribution by permutation symmetry. We call $\Omega$ the
total frequency given by the pump, so that the probe enters at $\omega-\Omega$; for the probe of
Eq.~\eqref{eq:delta_probe}, $\Ep(\omega-\Omega) = \Ep(\omega)\,e^{-i\Omega\tp}$. The term of first order gives
$\chi^{(1)}$, and the pump induced change of Eq.~\eqref{eq:dchi} is $\dchi = \dchi^{(2)}+\dchi^{(3)}+\dots$, with
\begin{align}
\dchi^{(2)}(\omega,\tp) &= 2\int d\Omega\,e^{-i\Omega\tp}\,\chi^{(2)}(\omega;\omega-\Omega,\Omega)\,\EP(\Omega) ,
\label{eq:dchi2}\\
\dchi^{(3)}(\omega,\tp) &= 3\int d\Omega\,e^{-i\Omega\tp}\int d\omega_1\,\EP(\omega_1)\,\EP(\Omega-\omega_1)
\nonumber\\
&\quad\times\chi^{(3)}(\omega;\omega-\Omega,\omega_1,\Omega-\omega_1) ,
\label{eq:dchi3}
\end{align}
and in general $\dchi^{(n+1)}$ contains $(n+1)\,\chi^{(n+1)}$ and $n$ pump fields. The term $\dchi^{(n+1)}$ describes a
probe photon at $\omega-\Omega$ that leaves the system at $\omega$ after the exchange of $n$ pump photons of total
frequency $\Omega$. No assumption is made on the pump. The frequency $\Omega$ enters only through the phase
$e^{-i\Omega\tp}$, which sets the dependence of $\dchi$ on the probe time.

% ====================================================================================================
\section{Pump and probe signal at the exciton peak of LiF}
\label{sec:lif}

% Plan: Outline.md, Sec. 3 (HSEX level only; computational details of the two approaches).

% ====================================================================================================
\section{Analysis and discussion of the non-equilibrium signal}
\label{sec:analysis}

% Plan: Outline.md, Sec. 4 (DC and AC parts, poles of chi3).

% ====================================================================================================
\section{Conclusions}
\label{sec:conclusions}

% Plan: Outline.md, Sec. 5.


\appendix
% ====================================================================================================
\section{Envelope of the pump}
\label{app:envelope}

% Derivation of the envelope at the probe time and of the accumulated envelope (Outline.md, Appendix A).

% ====================================================================================================
\section{Non-linear susceptibilities from the real-time polarization}
\label{app:extraction}

% Sec. 2 of the technical notes, shortened (Outline.md, Appendix B).


\bibliography{biblio}

\end{document}
